\ifdefined\KIOInclude\else
\documentclass[a4paper,12pt]{article}
\usepackage[utf8]{inputenc}
\usepackage[margin=2.5cm]{geometry}
\usepackage{amsmath,amssymb}
\begin{document}
\fi

\section*{Question 4: Mama Fua (25 marks)}

Mama Fua runs a laundry service in Nairobi. A customer brings \(n\) pieces
of clothing, and Mama Fua offers two ways to charge:

\begin{itemize}
    \item \(p\) shillings for each piece of clothing; or
    \item \(r\) shillings for each kilogram of clothing.
\end{itemize}

The customer may choose some pieces to be charged by weight. These pieces
are placed on the scale together, and their total weight is rounded
\emph{up} to the next whole kilogram. Every remaining piece is charged
separately at \(p\) shillings per piece.

For example, if the pieces charged by weight have a total weight of
\(2.01\) kilograms, the customer pays for \(3\) kilograms. If no pieces are
charged by weight, the weight charge is zero.

Find the minimum possible total bill.

\subsection*{Input}

The first line contains three integers \(n\), \(p\), and \(r\): the number
of pieces, the price per piece, and the price per kilogram.

The second line contains \(n\) weights \(w_1,w_2,\ldots,w_n\). Every weight
is positive and is written with exactly two digits after the decimal point.
Weights are given in kilograms.

\subsection*{Output}

Output one integer: the minimum possible bill in shillings.

\subsection*{Constraints}

\[
1\le n\le2\cdot10^5
\]
\[
10\le p,r\le10000
\]
\[
0.01\le w_i\le10000.00
\]

\subsection*{Example 1}

\paragraph{Input}
\begin{verbatim}
6 200 150
0.20 0.30 0.40 2.10 3.00 10.00
\end{verbatim}

\paragraph{Output}
\begin{verbatim}
750
\end{verbatim}

\subsection*{Example 2}

\paragraph{Input}
\begin{verbatim}
4 200 150
0.25 0.25 0.25 0.25
\end{verbatim}

\paragraph{Output}
\begin{verbatim}
150
\end{verbatim}

\subsection*{Explanation}

In the first example, the customer weighs the three lightest pieces
together. They weigh \(0.90\) kilograms, which is rounded up to \(1\)
kilogram and costs \(150\) shillings. The other three pieces cost
\(3\cdot200=600\) shillings, giving a total of \(750\).

In the second example, all four pieces are weighed together. Their total
weight is exactly \(1\) kilogram, so the bill is \(150\) shillings.

\medskip
\noindent\textbf{Mark Distribution [25 marks total]:}
\begin{itemize}
    \item \textbf{Algorithm (10 marks):} Write a clear, step-by-step algorithm or pseudocode in your answer booklet describing your solution.
    \item \textbf{Code Implementation (15 marks):} Implement your solution in any programming language of your choice (C, C++, Python, etc.) that passes all test cases within the time and memory limits.
\end{itemize}

\ifdefined\KIOInclude\else\end{document}\fi
